\cotefile{75}
\scanpage{26}
c) Orientation d'un ensemble fini $I$\inserted{(cardinal $n$).} La donnée de $I$ équivaut à celle d'un torseur
\[
 P=\operatorname{Isom}(\deleted{\unsure{[1,n]_{\mathbb Z}}}I_n,I)=\operatorname{numérot}(I)
\]
$(I_n\mathrel{\unsure{=|:}}\{m\in\mathbb N^{*}\mid m\leq n\}=\{\deleted{\unsure{0}}1,2,\ldots,n\})$ sous le groupe
\[
 \deleted{\unsure{\mathfrak S|S}}\mathfrak S_n=\mathfrak S_{I_n}=\operatorname{Aut}(I_n).
\]

On a un hom canonique ``signature''
\diagram{\begin{tikzcd}[ampersand replacement=\&,column sep=large]
\mathfrak S_n \arrow[r,"s=\mathrm{sgn}"] \& \{\pm1\} \& \text{noyau }\mathfrak A_n
\end{tikzcd}}
et par définition l'ens des 2 orientations de $I$ est
\[
 \omega_I=\deleted{\unsure{\operatorname{Isom}}(}P\times_{\mathfrak S_n}\{\pm1\}\simeq P/\mathfrak A_n .
\]

\uline{Compatibilités avec} a) \uline{et} b)

\uline{Avec} b) On a sur $E_{(I,k)}=k^{I}$ un produit scalaire canonique
\[
 (x,y)=\sum_{i\in I}x_i y_i
\]
pour lequel $(e_i)$ est une base orthonormale. Ainsi $\deleted{E}I$ \deleted{\unsure{donne}}\inserted{\unsure{donne}} naissance à un module quadratique $E(I,\unsure{x|k|X})$. Ceci posé, on \inserted{définit} une bijection canonique
\[
 \omega_{E(I,x)}\simeq\omega_I
\]
\margin{Si $i_1,\mathinner{\cdotp\cdotp}i_n$ est une numérotation de $I$, $e_{i_1}\mathinner{\cdotp\cdotp}e_{i_n}$ est une base \inserted{\unsure{ord.}} orthon. de $E$, \ldots}
par transitivité dans
\diagram{\begin{tikzcd}[ampersand replacement=\&,column sep=large]
 \mathfrak S_n \arrow[r,hook,"\mathrm{incl.}"] \arrow[rr,bend right=24,"\mathrm{sgn}"] \& \mathrm{O}(n,k) \arrow[r,"\det"] \& \{\pm1\}
\end{tikzcd}}

En particulier, on trouve (comp. avec a).),
\[
 \omega_{E(I,\mathbb R)\tikz[baseline=0pt]{\draw[->] (0,0) .. controls (.17,-.10) and (.31,-.22) .. (.49,-.22);}}\simeq\omega_I
\]
% The arrow starts at the closing parenthesis of E(I,R) and points down/right to the two-line note at the foot.
% The lower strokes of the second line meet the bottom edge of the scan.
\margin{considéré\\comme \unsure{vect.|variété} \unsure{réel|via|vide} \unsure{se pense}}
% archive sheet number: 25
